\scanpage{6}
\mspage{\unsure{291bis}}

\uline{38}\quad Explicit \unsure{determination} of the homomorphism
$\Pi'_{0,3}\xrightarrow{\sim}\mathfrak S_{0,3}^{+}$
(and of $\unsure{\Pi_{0,3}^{\prime\sim}}$
(unoriented triangulated cartographic group)
$\xrightarrow{\sim}\mathfrak S_{0,3}$).

I feel like coming back to the proposition of the preceding \S\
(page\ 262) with \inserted{\unsure{connected}} surface $X$, discrete group $G$
operating properly, the stabilizers \inserted{$G_x$} of the $x\in X$
\deleted{\unsure{being}} preserving the orientation and operating faithfully
in a neighborhood of $x$ \ldots\ taking $Y=X/G$, \unsure{equipped|equipped} with the
``signature'' $\underline d=d(X/Y)$ defined by the \unsure{ramification}
--- whence an \uline{\unsure{Isomorphism|isomorphism}}
\[
(1)\qquad
\Pi=\Pi_1((X,G),x)\xrightarrow{\sim}\Pi_1((Y,\underline d),y)=\Pi'
\]
($x\in X$, $y$ image of $x$ in $Y$)
which I would like to \unsure{study} more closely.

We have in $\Pi_1((Y,\underline d),y)$, for every
$b\in\unsure{\operatorname{Supp}}\underline d=Y^!$,
``\unsure{truncated-loop} subgroups'', more precisely
a conjugacy \unsure{class} \unsure{of|of the} \unsure{subgroups}
quotients of $\mathbb Z/d_b\mathbb Z$, and more precisely still,
a canonical outer homomorphism
\[
(2)\qquad
\unsure{T_b|T_y}\otimes_{\mathbb Z}\mathbb Z/d_b\mathbb Z
\longrightarrow\Pi'
\]
obtained by passage to the quotient from the outer homomorphism
\[
(3)\qquad
\unsure{T_b|T_y}\xrightarrow{\unsure{k_y|k_b}}
\Pi_1(Y\setminus Y^!)
\qquad(\longrightarrow\Pi_1((Y,\underline d)))
\]
(where $T_b$ is the free $\mathbb Z$-module of rank $1$ of
``twisted integers'' at $b$ on $Y$). On the other hand,
if $a\in X$ lies above $b$, the stabilizer $G_a$ is
\deleted{\unsure{thus}} a cyclic group of finite order $d_a=d_b$
--- We have \unsure{$=$} an obvious canonical isomorphism
\[
(4)\qquad
G_a\simeq T_a\otimes_{\mathbb Z}\mathbb Z/d_a\mathbb Z
\simeq T_b\otimes\mathbb Z/d_b\mathbb Z.
\]
\margin{(N.B. $T_a\simeq T_b$ \unsure{canonical\textsuperscript{ly}})}
\margin{5}

\scanpage{7}

On the other hand, to $a$ is \unsure{associated|associated} a conjugacy
class of \unsure{splittings} of the extension
\[
(5)\qquad
1\longrightarrow\Pi_1(X,x)\longrightarrow\Pi_1((X,G),x)
\longrightarrow G\longrightarrow1,
\]
(cf. \S\ 15)
\[
(6)\qquad
k_a:G_a\longrightarrow\Pi_1((X,G),x)
\]
these classes \inserted{of $\Pi_1(X,\unsure{x|e})$-conjugacy}
of \unsure{splittings} (if $X\ne S^2$) correspond
one-to-one to the fixed points of $G_a$ operating on $X$.
If we take the \uline{outer} \uline{homomorphism} \unsure{defined|defined}
by (6) (thus \unsure{conjugating|conjugating} not only by the
$\in\Pi_1(X)$, but by those of $\Pi_1(X,G)$ in its entirety)
we find a $1$-$1$ correspondence between \inserted{the set}
of outer homomorphisms (6) which lift the inclusion
\[
G_a\hookrightarrow G
\]
(\unsure{considered} as an outer homomorphism), and the set of orbits of $G$
in $X^!$, i.e.\ \deleted{between} the points of $Y^!$. In \unsure{summary}

\uline{Proposition}\quad We have a bijection between $Y^!$
and the set of ``maximal'' outer homomorphisms
\inserted{$\varphi=\varphi_b:\Gamma\longrightarrow\unsure{\Pi|\pi}$}
of \unsure{subgroups} $\Gamma\subset G$, \deleted{\unsure{(cyclic|(finite}}
$\Gamma\ne1$, into $\Pi_1(X,G)$ which lift the inclusion
\inserted{$\Gamma\subset G$} (the said $\Gamma$ are necessarily
cyclic, the $\Gamma$
\margin{6}

\scanpage{8}
\mspage{292bis}

in question \unsure{corresponding|corresponding} to a $b\in Y^!$
being the stabilizers $G_a$ in $G$ of the
\deleted{\unsure{$a\in(X^!)_a$|$a\in(X^!)_b$}}
$a\in X_b=X_b^!$. Moreover, the composite of the homomorphisms
\[
(7)\qquad
T_b\otimes\mathbb Z/n\mathbb Z
\overset{(4)}{\simeq}G_a
\xrightarrow{\varphi_a=\varphi_b}\Pi_1(X,G)
\overset{(1)}{\xrightarrow{\sim}}\Pi_1((Y,\underline d))
\]
is none other (we might have guessed!) than \unsure{``}the
``\unsure{truncated loops}'' homomorphism (2).

Let us now suppose that $X$ is of the form
$\widehat X\setminus S$, $\widehat X$ a compact surface,
$S$ a finite subset of $\widehat X$ \deleted{\unsure{thus}}
\[
(8)\qquad
\left\{\begin{aligned}
X&=\widehat X\setminus S\quad\text{thus}\quad
Y=\widehat Y\setminus T\quad\unsure{\ne\varnothing}\quad\text{\unsure{with}}\\
\widehat Y&=\widehat X/G
\end{aligned}\right.
\]

Let us consider\unsure{,} then, for every $\unsure{a|b}\in S$, the outer homomorphism
\[
(9)\qquad
k_a:T_a\longrightarrow\Pi_1(X)\hookrightarrow\Pi_1(X,G)
\]
which I am going to compare with the composite
\[
(10)\qquad
\deleted{\unsure{k_b:}}\quad T_b\xrightarrow{k_b}\Pi_1(Y)
\longrightarrow\Pi_1(Y,\underline d)
\]
(for $b=\widehat f(a)$), by means of the canonical isomorphism
\[
\deleted{\unsure{(12)}}\ (11)\qquad T_a\simeq T_b
\]
If $d_a$ again denotes the ramification index at $a$, we shall have
a commutative compatibility diagram:
\margin{7}

\scanpage{9}

\[
(12)\qquad
\diagram{\begin{tikzcd}[ampersand replacement=\&,column sep=large,row sep=large]
T_a \arrow[r,"k_a"] \arrow[d,"\wr\ \mathrm{can}"'] \& \Pi_1(X) \arrow[r] \& \Pi_1(X,G) \arrow[dd,"\wr"] \\
T_b \arrow[d,"d_a\cdot\mathrm{id}"'] \& \& \\
T_b \arrow[rr,"k_b"] \& \& \Pi_1(Y,\underline{d})
\end{tikzcd}}
\]

We are going to apply this to the \unsure{situation} of
$X=U_{0,3}^{\mathrm{an}}$, $G=\mathfrak{S}_3$,
$\widehat{X}=\mathbb{P}^{1\mathrm{an}}$, $\widehat{Y}\simeq\mathbb{P}^{1\mathrm{an}}$,
\unsure{with} the map of \unsure{passage} to the quotient
\unsure{given|given} by (page 286' (58))
\[
(13)\qquad
f(z)=\frac{3^3}{2^2}\frac{z^2(z-1)^2}{(z^2-z+1)^3}
\]
\[
f:\quad
\left\{\begin{array}{l}0\\1\\\infty\end{array}\right.\longrightarrow 0
\qquad
\left\{\begin{array}{l}2\\-1\\1/2\end{array}\right.\longrightarrow 1
\qquad
\left\{\begin{array}{l}-\overline{j}\\-j\end{array}\right.\longrightarrow\infty
\]

Thus we see, with the notations of \hspace{3em}, that the generators
$l'_0,l'_1,l'_\infty$ of $\Pi'_{0,3}$
(\unsure{satisfying}
\[
(14)\qquad l'_\infty l'_1 l'_0=1
\qquad {l'_1}^{2}={l'_\infty}^{3}=\unsure{\cdot}1\qquad )
\]
\unsure{correspond}, respectively, under the isomorphism
\[
(15)\quad
\begin{array}{ccl}
\mathfrak{S}_{0,3}^{+}
=\Pi_1(U_{0,3}^{\mathrm{an}},\mathfrak{S}_3)
&\xrightarrow{\sim}&\Pi'_{0,3}=\mathcal{T}_{1,1}^{+\prime}\\
\begin{array}{c}
\text{cartographic group}\\
\text{triangulated \unsure{weighted} oriented}\\
\text{\uline{mixed}}
\end{array}
&&\begin{array}{c}
\text{cartographic group \unsure{triangulated|triangul-}}\\
\text{oriented}\qquad \simeq\mathrm{Sl}(2,\mathbb{Z})'
\end{array}
\end{array}
\]
to the elements denoted
\deleted{$\unsure{\rho,\varepsilon_0,\varepsilon_1}$ \illegible{2}}
\[
(16)\qquad
\left\{\begin{array}{lll}
\varepsilon_0,\varepsilon_1,\varepsilon_\infty&\longmapsto&l'_0\\
\text{(such that $\varepsilon_i^2=l_i$)}&&\\
\sigma_0,\sigma_1,\sigma_\infty&\longmapsto&l'_1\\
\rho&\longmapsto&l'_\infty
\end{array}\right|\quad
\begin{array}{l}\mathrm{mod}\\\text{\unsure{conjugacy}}\end{array}
\]
\margin{or $\rho^{-1}$ ?}
\margin{8}

\scanpage{10}
\mspage{293}

\deleted{\unsure{The isomorphism}}

This shows that we can find
\inserted{in $\Pi=\mathfrak{S}_{0,3}^{+}$}
conjugates $\varepsilon'_0$ of $\varepsilon_0$, $\sigma'_0$ of $\sigma_0$,
such that we have
\[
(17)\qquad
\rho\sigma'_0\varepsilon'_0=1,\quad
\text{in addition to}\quad \rho^3={\sigma'_0}^{2}=1
\]
and in such a way that these relations \unsure{provide}
a \uline{presentation} of $\Pi$ (thus in particular,
that $\varepsilon'_0,\sigma'_0,\rho$ \unsure{generate} \unsure{$\}$} $\Pi$).
Drawing inspiration, moreover, from the isomorphism \deleted{more} which
\unsure{extends} (15) and is obtained in \unsure{like manner}
\[
(16)\qquad
\mathfrak{S}_{0,3}
\simeq\Pi_1\left(U_{0,3}^{\mathrm{an}},
\underbrace{\mathfrak{S}_3\times\unsure{\mathbb{Z}/2|\mathbb{Z}_2}}_{\mathcal{T}_{0,3}}\right)
\longrightarrow
\begin{array}{c}
\mathcal{T}'_{1,1}\\
\text{cartographic group}\\
\text{triangulated non-}\\
\text{oriented}\\
\simeq\mathrm{Gl}(2,\mathbb{Z})
\end{array}
\]
we should find on the first \unsure{side}
three \unsure{generators|generators} $\tau'_0,\tau'_1,\tau'_\infty$
(\unsure{corresponding} to the \unsure{analogous} generators
$\sigma_0,\sigma_1,\sigma_\infty$ of the
\unsure{corresponding|corresponding} cartographic groups) \unsure{satisfying|satisfying}
\[
(17)\qquad
{\tau'_0}^{2}={\tau'_1}^{2}={\tau'_\infty}^{2}=1,\quad
\left(\underbrace{\tau'_{\unsure{0|\infty}}\tau'_\infty}_{\begin{array}{c}\text{\unsure{corresponding} to}\\l'_1\end{array}}\right)^2
=\left(\underbrace{\tau'_1\tau'_0}_{\begin{array}{c}\text{corresponding to}\\l'_\infty\end{array}}\right)^3=1
\]
in such a way that these are \deleted{\unsure{some}} some
defining relations for $\mathfrak{S}_{0,3}$.
It would be necessary that
\deleted{$\tau'_\infty\tau'_0$ be a conjugate of $\sigma_0$}
\inserted{we have}
\[
(18)\qquad
\begin{array}{rclcrclcrcl}
\tau'_\infty\tau'_1&=&\varepsilon'_0;&\quad&
\tau'_0\tau'_\infty&=&\sigma'_0,&\quad&
\tau'_1\tau'_0&=&\rho\\
&&\updownarrow&&&&\updownarrow&&&&\updownarrow\\
&&l'_0&&&&l'_1&&&&l'_\infty
\end{array}
\]
\margin{9}